\section{Introduction}
\label{sec:introduction}

% Author text.

This paper sits at the intersection of several branches of the literature.

First, it relates to the literature on housing costs and fertility. A large
empirical literature finds that housing conditions shape childbearing.
Fertility is lower where rent per room is higher
\citep{simonTamuraHigherRents2009}, dwelling size predicts first births
\citep{stromHousingFirstBirths2010}, and higher house prices raise births
among owners but lower them among renters
\citep{dettlingHousePricesBirth2014,lovenheimFamilyWealthShocks2013}. Access
to mortgage credit also raises births
\citep{hacamoBabiesMortgageDeregulation2021,cummingMonetaryPolicyBirth2024,dettlingKearneyModernMortgage2025,vandoornikFazioRamadoraiSkrastinsHousingFertility2025}.
Two recent papers model housing and fertility jointly.
\citet{khorunzhinaMillerAmericanDream2022} estimate a dynamic model of female
labor supply, fertility and the first home purchase at given house prices,
and \citet{couillardBuildBabyBuild2025} estimates a model of fertility,
living arrangements and location in which rents are the cost of housing.
Unlike these papers, I study fertility together with saving, tenure and the
choice of space under uninsurable earnings risk, in a housing market where
the price is determined in equilibrium. This allows me to separate the
recurring cost of family space from the equity needed to buy it, and to
compare housing supply and mortgage credit within the same economy.

Second, the paper relates to work on earnings risk and fertility. Because a
child is an irreversible commitment, prudent parents who face more risk want
fewer children \citep{fraserIncomeRisk2001}. Quantitative life-cycle models
find that uninsurable earnings risk delays and reduces births
\citep{sommerFertilityChoiceLife2016,ejrnaesJorgensenFamilyPlanning2020,nakakuniFertilityIncomeUncertainty2026},
and \citet{deLaCroixPommeretChildbearing2021} show that risk gives the
postponement of a first birth an option value. Quasi-experimental studies
find that job displacement and job insecurity lower fertility
\citep{delBonoWeberWinterEbmer2012,clarkLepinteurJobInsecurity2022}. None of
these papers includes housing. In my model, part of the commitment a child
entails is the space it needs, paid for every period at a price set in the
housing market. The response of fertility to earnings risk therefore depends
on housing costs and on the down payment required to own family space.

Third, the paper builds on quantitative models of housing over the life
cycle. These models study tenure choice under earnings risk, down payments
and transaction costs
\citep{fisherGervaisHomeOwnership2011,halketVasudevSaving2014,boarGoreaMidriganLiquidity2022},
and the joint determination of house prices, rents and ownership
\citep{kaplanMitmanViolanteHousing2020,greenwaldGurenCredit2025}. In this
literature, rising earnings risk accounts for much of the decline in
homeownership among young households
\citep{fisherGervaisHomeOwnership2011,pazPardoHomeownership2024}. Household
size, however, is exogenous or absent. I make fertility a choice. A birth
raises the household's demand for space over its life cycle and, a
generation later, the number of households that compete for housing.
Housing policy therefore shapes both family size and the future demand for
housing.
